\documentclass[sigconf]{acmart}

\settopmatter{printacmref=false}
\renewcommand\footnotetextcopyrightpermission[1]{}
\pagestyle{plain}

%Course information
\acmConference
[CSE 6275/CSE 6391: AHCI, Assignment 2]
{}
{CSE Dept.}
{IUT}

\usepackage{graphicx}
\usepackage{booktabs}
\usepackage{multirow}
\usepackage{enumitem}
\usepackage{array}
\usepackage{hyperref}

\title{[Project Title]: Prototype Design and Construction for a Human-Centered AI System}

\author{Author One}
\affiliation{
  \institution{Department of Computer Science and Engineering\\
  Islamic University of Technology (IUT), OIC}
  \city{Dhaka}
  \country{Bangladesh}
}
\email{author1@iut-dhaka.edu}

\author{Author Two}
\affiliation{
  \institution{Department of Computer Science and Engineering\\
  Islamic University of Technology (IUT), OIC}
  \city{Dhaka}
  \country{Bangladesh}
}
\email{author2@iut-dhaka.edu}

\author{Author Three}
\affiliation{
  \institution{Department of Computer Science and Engineering\\
  Islamic University of Technology (IUT), OIC}
  \city{Dhaka}
  \country{Bangladesh}
}
\email{author3@iut-dhaka.edu}

\author{Hasan Mahmud}
\affiliation{
  \institution{SSL, Department of Computer Science and Engineering\\
  Islamic University of Technology (IUT), OIC}
  \city{Dhaka}
  \country{Bangladesh}
}
\email{hasan@iut-dhaka.edu}

\authornote{Corresponding author.}

\begin{document}

\begin{abstract}

Provide a concise summary of the prototype design (150--200 words). The abstract should briefly describe the problem domain, the purpose of the prototype, the type of prototype developed (e.g., low-fidelity or high-fidelity), and the tools used to build it. Emphasize the Human-Centered Artificial Intelligence (HCAI) considerations addressed in the design such as transparency, explainability, trust, and meaningful human control \cite{shneiderman2022hcai}. The abstract should summarize the overall design contribution and explain how this prototype supports the development of the proposed HCAI system.

\end{abstract}

\keywords{Human-Centered AI, Prototyping, Interaction Design, Human-AI Interaction}

\maketitle


\section{Introduction}

\subsection{Background and Motivation}

Introduce the problem domain and describe the real-world problem that your proposed system aims to address. Explain why AI-based assistance may help solve this problem and why it is important to design such systems using human-centered design principles.

Human-Centered Artificial Intelligence (HCAI) promotes the development of AI systems that augment human capabilities while preserving human control, accountability, and transparency \cite{shneiderman2022hcai}. Designing AI-assisted systems therefore requires careful consideration of usability, trust, and human oversight.

\subsection{From Requirements to Design}

Briefly summarize the key outcomes from Assignment 1, including the stakeholders identified, the major user needs, and the system requirements. These requirements should guide the design of the prototype presented in this assignment.

Interaction design literature emphasizes that understanding users and their context is a crucial step before developing interactive systems \cite{preece2015interaction,sharp2019interaction}.

\subsection{Objectives and Contributions}

The objectives of this assignment are to translate previously identified requirements into a concrete system design and to construct a prototype that demonstrates human--AI interaction. The prototype should illustrate key system functionality, interface layouts, and interaction workflows that support user goals.

\section{Related Work and Design Inspirations}

Human–Computer Interaction research highlights the importance of iterative design and prototyping when developing interactive systems. Prototypes help designers explore alternative interaction strategies and test usability early in the development process \cite{sharp2019interaction}.

In addition, goal-directed design approaches recommend grounding system design in personas, scenarios, and behavioral goals derived from empirical user data \cite{cooper2014aboutface}. These approaches are particularly important when designing AI-based systems, where inappropriate automation or opaque decision-making may undermine user trust.

Recent work in Human-Centered AI also stresses the need for transparency, explainability, and meaningful human oversight in AI-assisted decision systems \cite{amershi2019guidelines}.

\section{Design Goals and HCAI Principles}

\subsection{Usability Goals}

Describe usability goals such as effectiveness, efficiency, learnability, and error prevention. These goals ensure that users can easily understand and interact with the system.

\subsection{User Experience Goals}

Explain the desired user experience when interacting with the system. Examples include fostering trust in AI recommendations, minimizing cognitive load, and providing clear feedback and guidance.

\subsection{Human-Centered AI Principles}

Explain how the prototype incorporates HCAI principles such as transparency, explainability, reliability, and human oversight. Designing systems that clearly communicate AI behavior helps users maintain confidence and control over automated decisions \cite{amershi2019guidelines}.

\section{Conceptual Design of the System}

\subsection{System Overview}

Provide a high-level overview of the system and its purpose. Describe the key functionalities and the intended user group.

\subsection{Conceptual Model}

Describe the conceptual model of the system and how users interact with different components. Conceptual models help users understand how a system works and how their actions affect system outcomes \cite{preece2015interaction}.

\subsection{Human--AI Interaction Model}

Explain how the AI component supports user tasks. For example, the AI system may generate recommendations, provide predictions, or assist users in making decisions.

\section{Prototype Design Methodology}

\subsection{Prototyping Approach}

Explain the design process used to construct the prototype, including brainstorming, sketching, wireframing, and iterative design. Prototyping allows designers to explore design alternatives and refine interaction strategies before full implementation \cite{sharp2019interaction}.

\subsection{Tools and Technologies}

Describe the tools used to create the prototype (e.g., Figma, Balsamiq, Adobe XD, HTML/CSS, Dialogflow).

\subsection{AI Behavior Simulation}

If the prototype includes AI functionality, describe how AI responses are simulated within the prototype environment.

\section{Prototype Artifacts and Interface Design}

\subsection{Wireframes and Sketches}

Present early design artifacts such as sketches or wireframes that illustrate the structure and navigation of the system interface.

\begin{figure}[h]
\centering
\includegraphics[width=0.8\linewidth]{wireframe.png}
\caption{Example wireframe of the proposed interface}
\end{figure}

\subsection{Interface Screens}

Describe the main interface screens and their functions. Explain design decisions related to layout, interaction elements, and information presentation.

\subsection{Storyboards and Interaction Scenarios}

Provide scenarios demonstrating how users interact with the system in realistic situations.

\section{User Interaction and Task Flow}

\subsection{Task Scenarios}

Describe typical tasks users perform when interacting with the system.

\subsection{Interaction Flow}

Explain the sequence of user actions required to complete key tasks.

\subsection{AI Interaction Flow}

Describe how the AI system provides recommendations, feedback, or explanations during interaction.

\section{Implementation Details}

\subsection{Tools and Platforms}

Explain the tools or frameworks used to construct the prototype.

\subsection{Prototype Features}

Describe the functionality implemented in the prototype.

\subsection{Constraints and Limitations}

Discuss limitations such as simulated AI behavior or incomplete system functionality.

\section{Discussion}

\subsection{Alignment with Requirements}

Explain how the prototype satisfies the requirements identified in Assignment 1.

\subsection{Design Trade-offs}

Discuss major design decisions and trade-offs made during prototype development.

\subsection{Human--AI Interaction Challenges}

Reflect on challenges encountered when designing AI-assisted interactions.
\section{Design Implications}
Discuss how the insights from your prototype design may inform future HCAI systems or interaction design practices.

\section{Limitations}
Discuss limitations of the prototype, design assumptions, or methodological constraints.

\section{Conclusion}

Summarize the prototype design and explain how it supports the development of the proposed Human-Centered AI system. Briefly discuss future work such as usability evaluation and iterative refinement.

\section*{Acknowledgments}

The writing sections and subsections given in the template are defined by the course teacher. The points to be written in the report are decided based on the related topics on prototype design and construction covered in the class. The orientation, writings, and template generation in Overleaf were refined with the help of AI. 

\bibliographystyle{ACM-Reference-Format}
\bibliography{references}

\end{document}